\documentclass[a4paper, 12pt]{amsart}
\usepackage[french]{babel}
\usepackage[utf8]{inputenc}
\usepackage{amssymb}
\usepackage{graphicx}
\usepackage{coursebook}
\fthmstyle{plain}
\newfancythm{fthm}{Théorème}
\fthmstyle{defn}
\newfancythm{fdefn}{Définition}
\fthmstyle{ex}
\newfancythm{fex}{Exercice}
%\DeclareMathOperator*{\ln}{ln}
\newcommand{\pd}[2]{\ensuremath{\frac{\partial #1}{\partial #2}}}
\usepackage{pgf,tikz}
\usetikzlibrary{arrows, shapes}
\usetikzlibrary[patterns]
\usetikzlibrary{calc,fadings,decorations.pathreplacing, bending} 
\usetikzlibrary{decorations.markings}
\usetikzlibrary{decorations.pathmorphing}
\usetikzlibrary{positioning}
\usetikzlibrary{intersections}
\usetikzlibrary{hobby}

\title{Questions de cours}
\setcounter{section}{0}
\graphicspath{{../images}}
\begin{document}
\maketitle
%\input{ex6.tex}
\section{Questions de cours}
Les questions proposées dans ce document portent sur le cours et sont regroupées en thèmes. Lors de l'examen final, seules les questions de cette liste pouront être
posées, avec un minium de une dans chaque thème.
\subsection{Thème holomorphie}
Dans cette partie, le terme domaine d'holomorphie désigne le plus grand ouvert
dans lequel une application est holomorphe.
\begin{enumerate}
\item Quel est le domaine d'holomorphie de $f : z \mapsto z^4$ ?
\item Quel est le domaine d'holomorphie de $f : z \mapsto |z|$ ?
\item)Quel est le domaine d'holomorphie de $f : z \mapsto 1/ \cos(z)$ ?
\item Quel est le domaine d'holomorphie de $f : z \mapsto 1/(1 + z^2 )$ ?
\item Quel est le domaine d'holomorphie de $f : z \mapsto e^z /(1 + z^2 )$ ?
\item L'application $z \mapsto z\overline{z}$ est-elle holomorphe en $z = 1$ ?
\item Expliquer pourquoi l'application $z \mapsto z^2$ est holomorphe dans $\C$ et donner
sa dérivée.
\item Si $f, g$ sont deux applications holomorphes dans $\C$ , l'application $f.g$ est-elle
holomorphe ? Quelle est sa dérivée ?
\item Si $f, g$ sont deux applications holomorphes dans $\C$ , l'application composée
$f \circ g$ est-elle holomorphe ? Quelle est sa dérivée ?
\item Si $f, g$ sont deux applications holomorphes dans $\C$, en quels points l'appli-
cation $f /g$ est-elle holomorphe ? Quelle est sa dérivée ?
\end{enumerate}
\subsection{Thème séries entières et fonctions analytiques}
\begin{enumerate}
\item Enoncer le critère de d'Alembert pour une série de terme général $a_n$.
\item A l'intérieur de son disque ouvert de convergence, une application définie
par une série entière est-elle holomorphe ? Si oui, comment exprime-t-on sa
dérivée ?
\item Quel est le rayon de convergence de la série $\sum_{n \geq 0} x^n / n!$ ?
\item Quel est le rayon de convergence de la série $\sum_{n \geq 0} (-1)^n x^n / n!$ ?
\item
\item Énoncer le théorème des zéros isolés.
\end{enumerate}
\subsection{Thème formule de Cauchy}
\begin{enumerate}
	\item Soit $\gamma$ le lacet $t \in [0, 2\pi] \mapsto e^{it}$ . Que vaut :
	\[
	\int_\gamma \frac{dz}{z} \quad ?
	\] 
	\item Soit $\gamma$ le lacet $t \in [0, 2\pi] \mapsto e^{it}$ . Que vaut :
	\[
	\int_\gamma \frac{\sin z}{z} dz \quad ?
	\] 
	\item Soit $\gamma$ un lacet simple de classe $C^1$ . Soit $f$ une application holomorphe dans
	$\C$ . Que vaut :
		\[
	\int_\gamma f(z) dz \quad ?
	\] 
	\item Soit  $\gamma$  un lacet simple de classe $C^1$ ne passant pas par $z_0 \in \C$ . Soit $f$ une
	application holomorphe dans C . Que vaut :
	\[
\int_\gamma \frac{f(z)}{z-z_0} dz \quad ?
\] 
		\item Soit  $\gamma$  un lacet simple de classe $C^1$ ne passant pas par $z_0 \in \C$ . Soit $f$ une
	application holomorphe dans C . Que vaut :
	\[
	\int_\gamma \frac{f(z)}{(z-z_0)^2} dz \quad ?
	\] 
\end{enumerate}
\subsection{Thème série de Laurent et résidus}
\begin{enumerate}
	\item Donner le développement en série de Laurent de $z \mapsto e^z$ au voisinage de 0.
	\item Donner le développement en série de Laurent de $z \mapsto 1+z+z^{-1}$ au voisinage de 0.
	\item Enoncer le théorème des résidus.
	\item Quel est le résidu en 0 de l'application $z \mapsto cos(z)/z$ ?
	\item Quel est le résidu en $i$ de l'application $z \mapsto 1/(1 + z^2 )$ ?
	\item Soit $\gamma$ le lacet $t \in [0, 2\pi] \mapsto 2e^{it}$ . Que vaut :
\[
\int_{\gamma} \frac{e^z}{1+z^2}dz\quad ?
\]
	\item Soit $\gamma$ le lacet $t \in [0, 2\pi] \mapsto 2e^{it}$ . Que vaut :
	\[
	\int_{\gamma} \frac{z}{1+z}dz\quad ?
	\]
\end{enumerate}
\end{document}